% experiments.tex
%
% Post-submission experiment record: crossing the bridge from the diagnostic recovery
% task to controllable generation.
%
% SCOPE NOTE. This document is a NEW addition under Doc/final/future_work/. It does not
% modify, include, or depend on anything in Doc/final/thesis/, Doc/final/proposal/ or
% Doc/final/beamer/, all three of which are unchanged from the submitted state. It is a
% standalone article and compiles on its own with pdflatex.
%
% The authoritative machine-readable record is CTG_EXPERIMENTS_LOG.md in the repository
% root, and every number below carries the JSON path it came from.

\documentclass[11pt,a4paper]{article}

\usepackage[T1]{fontenc}
\usepackage[utf8]{inputenc}
\usepackage[margin=2.5cm]{geometry}
\usepackage{booktabs}
\usepackage{longtable}
\usepackage{amsmath}
\usepackage{amssymb}
\usepackage{bm}
\usepackage{graphicx}
\usepackage{xcolor}
\usepackage{listings}
\usepackage[hidelinks]{hyperref}

\lstset{basicstyle=\ttfamily\footnotesize, breaklines=true, frame=single,
        columns=fullflexible, xleftmargin=0.5em}

\newcommand{\ve}{\bm{e}}
\newcommand{\vs}{\bm{s}}
\newcommand{\vx}{\bm{x}}
\newcommand{\vg}{\bm{g}}
\newcommand{\vh}{\bm{h}}
\newcommand{\file}[1]{\texttt{\small #1}}

\title{Post-Submission Experiments:\\
From the Diagnostic Task to Controllable Generation}
\author{Experiment record, generated during defense preparation}
\date{2026-07-29}

\begin{document}
\maketitle

\begin{abstract}
The submitted thesis diagnoses why an input-embedding gradient of a frozen
autoregressive likelihood carries no usable directional signal, and demonstrates a
repair, the relaxed token-indicator (one-hot) surrogate, on a single-token recovery
task. It never crosses the bridge to controllable generation. This document records the
experiments run after submission to cross that bridge, or to measure exactly why it
cannot be crossed. Four findings are load-bearing. First, the repair does not need a
gradient at all: the self term $\log p(v \mid \vx_{<i})$, which the forward pass already
computes, matches the full token-indicator surrogate at one quarter of the cost, with
zero backward passes. Second, the thesis's headline $40\%$ recovery figure reproduces
exactly at its own $n=50$ but settles at $26.5\%$ at $n=200$, which removes the claimed
parity with a purpose-trained diffusion model while leaving the zero-versus-nonzero
finding untouched. Third, implementing MuCoLa's actual energy in the same code makes the
continuous sampler produce fluent text where the thesis's energy produces word salad,
which confirms the scope restriction the submitted thesis already states. Fourth, a
likelihood-ranked shortlist and a sentiment constraint want almost unrelated tokens, and
the classifier's own first-order expansion has a rank correlation near zero with the
change it is expanding, which supplies a mechanism for the weakness of gradient-based
constraint guidance in this setting.
\end{abstract}

\tableofcontents

\section{Purpose, scope and what is authoritative}

The thesis is submitted and nothing in it is changed by this work. This document exists
so that the author can answer, with numbers, the question the thesis leaves open: is the
sampler actually useful for controllable generation?

Four questions were posed, in priority order.

\begin{description}
\item[Q1 Cost.] The output-side surrogate decomposes into a \emph{self} term,
  $\log p(v \mid \vx_{<i})$, which the energy's own forward pass already computes and
  which is therefore free, and a \emph{future} term describing how substituting $v$
  changes everything downstream, which costs either one backward pass (linearized) or
  $k$ forward passes (exact). Is the future term worth its cost?
\item[Q2 Conditioning.] Do bidirectional proposals, which see right context, beat the
  autoregressive output-side proposal when the controlled attribute is global?
\item[Q3 Mechanism.] How should a constraint enter: exact evaluation on a shortlist, a
  first-order surrogate of the classifier, both, or only the accept step?
\item[Q4 Usefulness.] At matched compute, does the exact-energy Metropolis--Hastings
  chain beat MuCoLa on the adherence-fluency trade-off, and does it produce more
  distinct valid solutions?
\end{description}

The authoritative machine-readable record is \file{CTG\_EXPERIMENTS\_LOG.md} in the
repository root. Every number in this document carries the path of the JSON it came
from. Suggested defense slides are in \file{SUGGESTED\_BACKUP\_SLIDES.md} and are
suggestions only; nothing has been inserted into the deck.

\subsection{Relationship to the submitted documents}

This file lives in \file{Doc/final/future\_work/} and is a new addition. It does not
modify, include or depend on \file{Doc/final/thesis/}, \file{Doc/final/proposal/} or
\file{Doc/final/beamer/}. Those three directories are byte-for-byte unchanged from the
submitted state; their checksums are recorded in the repository log.

\section{Part 0: audit of what already existed}

\subsection{The constrained task was reused, not rebuilt}

\file{scripts/run\_constrained.py} already defines the whole constrained-generation
task, and it was imported rather than reimplemented: the fifteen PPLM discriminator
prompts verbatim from MuCoLa's repository, the twenty-token continuation span, the
centroid and random-token initializations, the two setups, the five constraint-mode
arms, and the sentiment head. Only the sampler and the recording are new.

\subsection{What \texttt{policy\_onehot} actually computes}

This determined the arm list and was established by reading the code rather than from
memory. In \file{core/dls.py} the proposal logit for candidate $v$ is built as
\[
  t_2[v] \;=\; \underbrace{\tfrac{1}{2}\,\vg^\top\!\big(\ve(v) - \vs\big)}_{\text{embedding-gradient (future) term}}
  \;+\; \underbrace{\tfrac{1}{2}\,\log p(v \mid \vx_{<m})}_{\text{self term, added only for \texttt{policy\_onehot}}}
\]
so \texttt{policy\_onehot} is the self term \emph{plus} the input-embedding-gradient
term, not the self term alone. It costs one forward, one backward and one extra no-grad
forward per step, and the same again for the Metropolis reverse evaluation.

Consequently the thesis's headline one-hot result had never been decomposed into its two
halves. Doing so is the first experiment below, and it is cheap.

\section{Study A: is the future term worth its cost?}

\subsection{Arms and pre-registration}

\begin{table}[htbp]
\centering
\begin{tabular}{llc}
\toprule
Arm & Proposal score for candidate $v$ & Model calls per step \\
\midrule
A1 \texttt{policy\_self}    & $\tfrac12 \log p(v \mid \vx_{<m})$ & 1 fwd, 0 bwd \\
A2 \texttt{policy\_onehot}  & A1 $+\ \tfrac12 \vg^\top(\ve(v)-\vs)$ & 2 fwd, 1 bwd \\
A3 \texttt{policy\_exact\_k}& $\tfrac12\big(E(v) - E(\text{cur})\big)$ on a shortlist & 1 fwd $+\ k$ fwd \\
A4 \texttt{policy}          & $\tfrac12 \vg^\top(\ve(v)-\vs)$ alone & 1 fwd, 1 bwd \\
A5 \texttt{uniform}         & constant & 1 fwd, 0 bwd \\
\bottomrule
\end{tabular}
\caption{Study A arms, all inside the exact-energy Metropolis--Hastings chain.}
\end{table}

Three predictions were recorded in advance: (1) A1 approximately equals A2 at strictly
lower cost; (2) A3 buys little over A1 relative to its cost; (3) A4 stays at the floor
with A5.

\subsection{Two framings}

\textbf{Framing 1, full vocabulary.} Every arm proposes over all $50{,}257$ tokens with
the Langevin distance term present, in the cell that produced the thesis's headline
one-hot result ($\epsilon: 10.5 \to 0.1$, $T=1.0$, Metropolis on, gradient normalization
off). Directly comparable to the published figures. A3 cannot appear here: exact scoring
of the full vocabulary would cost $50{,}257$ forward passes per step.

\textbf{Framing 2, matched shortlist.} Every arm proposes over the same frozen top-$k$
candidate set ($k \in \{16,64,256\}$) scored only by its own surrogate, distance term
dropped. This is the clean comparison, because the exact arm can only ever see a
shortlist, so the restriction is charged to every arm equally.

\subsection{Three defects found before any real run}

Each of these would have produced a confident and wrong answer, and each was caught by a
smoke test.

\begin{enumerate}
\item \textbf{Shortlist excludes the incumbent.} The exact arm put zero mass on the token
  the chain currently held, so the reverse probability was zero, $\log\alpha$ was
  $-\infty$, and acceptance was $0.0\%$ on every sequence.
\item \textbf{The distance term pins a shortlisted chain.} With $k=64$ and the Langevin
  term $t_1$ present, all five arms returned an identical final KL of $6.777$ and zero
  exact matches, because $t_1$ scores the incumbent at exactly $0$ and every rival at
  $-\lVert \ve(v)-\vs \rVert^2 / 2\epsilon$. Over $50$k candidates that is one preference
  among many; over $65$ it is decisive.
\item \textbf{A state-dependent shortlist is not reversible.} A shortlist recomputed at
  each state as ``top-$k$ here plus wherever I am'' gives the reverse move probability
  zero as soon as the chain leaves its starting token, which it does immediately.
  Acceptance was again $0.0\%$ while the proposal itself was healthy: measured directly,
  $24.6\%$ of the proposal mass sat on the ground-truth token and the move was worth
  $+19.8$ nats, and it was refused because the reverse term was $-\infty$. Fixed by
  freezing the support per sequence, which makes each arm an independence sampler on a
  fixed candidate set, a valid Metropolis--Hastings kernel.
\end{enumerate}

A regression gate was run after every sampler edit: the flagship archived configuration
reproduces \file{results/grid/gpt2\_v2/} bit-identically, maximum absolute difference
$0.000\mathrm{e}{+}00$ on all three recorded channels.

\subsection{Result, framing 1 (full vocabulary), $n=200$}

Source: \file{results/ctg/studyA/A.fullvocab.*.json}, index
\file{results/ctg/studyA\_summary.json}.

\begin{table}[htbp]
\centering
\begin{tabular}{lrrrrrr}
\toprule
Arm & Exact \% & Final KL & Accept \% & fwd-eq/seq & GPU s/seq & GPU s/accept \\
\midrule
A1 \texttt{policy\_self}   & 24.5 & 4.598 & 25.2 & \textbf{100} & \textbf{4.18} & \textbf{0.332} \\
A2 \texttt{policy\_onehot} & 26.5 & 4.584 & 24.3 & 400 & 10.71 & 0.881 \\
A4 \texttt{policy}         &  0.0 & 6.243 & 11.1 & 300 &  7.71 & 1.391 \\
A5 \texttt{uniform}        &  0.0 & 6.671 &  9.9 & 100 &  4.24 & 0.857 \\
\bottomrule
\end{tabular}
\caption{Study A, full vocabulary. Paired bootstrap contrasts against A1:
A2 gives $\Delta$KL $-0.015$ $[-0.361,+0.331]$ and $\Delta$exact $+2.0$ points
$[-0.5,+5.0]$ at four times the forward-equivalent cost; A4 gives $-24.5$ points
$[-30.5,-18.5]$; A5 gives $-24.5$ points $[-30.5,-18.5]$.}
\end{table}

Prediction 1 is confirmed strongly. The entire linearized future term, the thing that
costs the backward pass, buys no measurable quality. Put the way a defense question
wants it: \emph{the thesis's recovery result is obtainable with no backward pass at
all}. \texttt{policy\_self} never calls autograd once, which the run's own counter
certifies, and it recovers $24.5\%$ where the input-embedding gradient recovers $0.0\%$
on the same two hundred sequences.

This also answers a question the submitted thesis explicitly leaves open. Section 5.6.1
observes that at the final position the future term vanishes so the token-indicator
derivative reduces exactly to $\log p(v \mid \vx_{<i})$, and adds that ``in the interior
the two come apart, because the future term is then nonzero''. Measured in the interior,
they do not come apart to any degree the data can resolve.

\subsection{Result, framing 2 (matched shortlist), $n=200$ per cell}

\begin{table}[htbp]
\centering
\begin{tabular}{llrrrrr}
\toprule
$k$ & Arm & Exact \% & Final KL & Accept \% & fwd-eq/seq & GPU s/seq \\
\midrule
16 & \texttt{policy\_self}    & 30.0 & 4.788 & 30.4 & 100 & 4.22 \\
16 & \texttt{policy\_onehot}  & 30.0 & 4.948 & 30.4 & 400 & 10.57 \\
16 & \texttt{policy\_exact\_k}& 30.5 & 4.879 & 80.2 & 117 & 4.41 \\
16 & \texttt{policy}          & 30.5 & 4.814 & 17.0 & 400 & 10.42 \\
16 & \texttt{uniform}         & 31.5 & 4.580 & 16.6 & 100 & 5.56 \\
\midrule
64 & \texttt{policy\_self}    & 29.5 & 4.706 & 25.9 & 100 & 5.44 \\
64 & \texttt{policy\_onehot}  & 28.0 & 4.822 & 25.7 & 400 & 10.63 \\
64 & \texttt{policy\_exact\_k}& \textbf{31.0} & \textbf{4.150} & 73.0 & 165 & \textbf{4.68} \\
64 & \texttt{policy}          & 21.5 & 4.716 & 12.6 & 400 & 11.66 \\
64 & \texttt{uniform}         & 18.5 & 4.875 & 13.4 & 100 & 4.26 \\
\midrule
256 & \texttt{policy\_self}    & 26.5 & 4.638 & 23.0 & 100 & 4.24 \\
256 & \texttt{policy\_onehot}  & 25.0 & 4.660 & 22.7 & 400 & 10.91 \\
256 & \texttt{policy\_exact\_k}& \textbf{31.0} & \textbf{4.037} & 66.6 & 357 & 5.58 \\
256 & \texttt{policy}          &  8.5 & 5.271 & 11.1 & 400 & 11.42 \\
256 & \texttt{uniform}         &  8.0 & 5.195 & 12.4 & 100 & 4.22 \\
\bottomrule
\end{tabular}
\caption{Study A, matched shortlist. All arms propose over the same frozen candidate
set.}
\end{table}

\begin{table}[htbp]
\centering
\begin{tabular}{llrrl}
\toprule
$k$ & Contrast & $\Delta$KL & $\Delta$exact (pts) & Cost \\
\midrule
16  & onehot $-$ self  & $+0.160\ [-0.011,+0.360]$ & $+0.0\ [-2.0,+2.0]$ & $\times 4.00$ fwd-eq \\
16  & exact$_k$ $-$ self & $+0.090\ [-0.264,+0.442]$ & $+0.5\ [-3.0,+4.0]$ & $\times 1.17$ fwd-eq \\
64  & onehot $-$ self  & $+0.116\ [-0.121,+0.363]$ & $-1.5\ [-4.0,+0.5]$ & $\times 4.00$ fwd-eq \\
64  & \textbf{exact$_k$ $-$ self} & $\bm{-0.555\ [-0.977,-0.161]}$ & $+1.5\ [-2.0,+5.5]$ & $\times1.65$ fwd-eq, $\bm{\times 0.86}$ wall \\
256 & onehot $-$ self  & $+0.023\ [-0.212,+0.269]$ & $-1.5\ [-4.0,+1.0]$ & $\times 4.00$ fwd-eq \\
256 & \textbf{exact$_k$ $-$ self} & $\bm{-0.601\ [-1.008,-0.205]}$ & $\bm{+4.5\ [+0.0,+9.0]}$ & $\times 3.57$ fwd-eq \\
\bottomrule
\end{tabular}
\caption{Paired contrasts against the self term, $10{,}000$-sample bootstrap, paired on
\texttt{sample\_idx}. The corruption is deterministic per index and identical across
arms, so the pairing needs no rerun.}
\end{table}

\paragraph{Prediction 2 is refuted, and is reported as refuted.}
The exact future term does buy a real improvement once the shortlist is wide enough for
ranking to matter: $-0.555$ nats at $k=64$ and $-0.601$ at $k=256$, both confidence
intervals excluding zero, plus $+4.5$ exact-match points at $k=256$. It is also cheap,
which was the unexpected part: at $k=64$ it costs $1.65$ times the forward-equivalents
and $0.86$ times the wall-clock, that is, it is faster in real time than the arm it
beats. Two measured reasons: its acceptance rate is $66$ to $80\%$ against $23$ to $30\%$
for every other arm, so it wastes far fewer steps; and with one masked position the $k$
candidate energies do not depend on the current token, so they are computed once per
sequence and reused. That second economy is task-specific and does not survive to a task
where many positions move.

\subsection{The support, not the scoring rule, is the dominant factor}

\begin{table}[htbp]
\centering
\begin{tabular}{lrrrr}
\toprule
Arm & Full vocab & $k=256$ & $k=64$ & $k=16$ \\
\midrule
\texttt{uniform} (no scoring at all) & 0.0 & 8.0 & 18.5 & \textbf{31.5} \\
\texttt{policy} (input-embedding gradient) & 0.0 & 8.5 & 21.5 & \textbf{30.5} \\
\texttt{policy\_self} & 24.5 & 26.5 & 29.5 & 30.0 \\
\midrule
gradient minus uniform (points) & 0.0 & $+0.5$ & $+3.0$ & $-1.0$ \\
\bottomrule
\end{tabular}
\caption{Exact recovery by shortlist width. At $k=16$ every arm lands between $30.0$ and
$31.5$ percent, including the arm that scores nothing.}
\end{table}

This does \emph{not} rehabilitate the input-embedding gradient. Its correct control is
the uniform arm, and it tracks that floor at every width, never separating from it. What
it does mean is that any claim of the form ``surrogate $X$ works in a shortlisted
sampler'' must be quoted against the uniform-on-the-same-shortlist floor rather than
against zero. That floor was not measured anywhere in the submitted thesis, and holding
the support fixed is a stronger statement of the null than the thesis makes, because it
rules out the objection that the null was an artifact of proposing over $50{,}257$
candidates.

\section{A correction to a claim the submitted thesis makes}
\label{sec:contradiction}

The submitted thesis, \file{chapters/05\_results.tex} line 413, states that the
token-indicator surrogate ``recovers $36$ to $40$ percent of corrupted tokens'', and
line 462 describes the outcome as ``recovery comparable to a purpose-trained diffusion
model''.

\begin{table}[htbp]
\centering
\begin{tabular}{lrl}
\toprule
$n$ & Exact recovery & $95\%$ CI \\
\midrule
50 (the sweep cell as published) & $40.0\%$ & $[26.4,\ 53.6]$ \\
200 (same configuration, rerun) & $\bm{26.5\%}$ & $[20.4,\ 32.6]$ \\
\bottomrule
\end{tabular}
\caption{The same cell at two sample sizes. Sources
\file{results/grid/rev3/ohsweep\_e10p5\_t1p0.json} and
\file{results/ctg/studyA/A.fullvocab.policy\_onehot.json}.}
\end{table}

The published number is reproducible and is not an error: rerunning the configuration
and taking the same first fifty sequences gives exactly $40.0$ percent. What the larger
sample shows is that $40.0$ percent was an $n=50$ estimate which regresses to $26.5$
percent at $n=200$, below the range the text states.

The comparability claim is the part that does not survive. The comparators it is set
against, SEDD-small $38.5\%$, SEDD-medium $39.0\%$ and RoBERTa-large $44.5\%$, were all
measured at $n=200$. The thesis therefore matches an $n=50$ estimate against $n=200$
estimates. At equal $n$ the token-indicator sampler reaches $26.5\%$ $[20.4, 32.6]$ and
is below both rather than comparable to either.

\paragraph{What survives unchanged, and it is the load-bearing part.}
The qualitative claim is not weakened and is now measured at four times the sample size:
the token-indicator surrogate recovers $26.5$ percent where the input-embedding gradient
recovers $0.0$ percent on the same sequences, a paired difference of $24.5$ points, CI
$[18.5, 30.5]$. The direction, the mechanism and the diagnosis are untouched. What
changes is the size of the number and the claim of parity with a diffusion model.

\section{The \texttt{mucola\_faithful} arm}

\subsection{Why it exists}

The thesis's continuous sampler is a variant of MuCoLa, not MuCoLa. Following Kumar,
Paria and Tsvetkov (2022), section 3, the MuCoLa language-model term is
\[
  \log P(\tilde{\ve}_{n+1} \mid \tilde{\ve}_{1:n})
  = \vh_n^\top \tilde{\ve}_{n+1} + b_{n+1}
    - \log \sum_j \exp\!\big(\vh_n^\top \ve_j + b_j\big),
\]
with the continuous state $\tilde{\ve}$ in the softmax \emph{numerator} and the real
embedding table in the denominator. GPT-2's output head is tied to the input embedding
and carries no bias, so $b=0$. The state therefore receives gradient by both paths the
paper names: directly, with gradient exactly $\vh_n$, and through the following hidden
states.

The repository's energy instead gathers the score at the projected token \emph{index},
so the direct path is absent by construction. That is the object the thesis's null
characterizes.

\subsection{Verification against the paper}

Run before the arm was used for anything.

\begin{table}[htbp]
\centering
\begin{tabular}{p{5.6cm}p{4.4cm}p{4.4cm}}
\toprule
Check & Expected & Measured \\
\midrule
MuCoLa energy at a real token sequence & equals the ordinary sequence log-likelihood, since $\tilde{\ve}=\ve(x)$ recovers the standard softmax & $-46.909576$ vs $-46.909576$, \textbf{abs.\ diff.\ $0.000\mathrm{e}{+}00$} \\
Gradient at the final position, MuCoLa energy & exactly $\vh_{n-1}$, the paper's direct path & norm $50.5592$ vs $\lVert \vh_{n-1}\rVert = 50.5592$, \textbf{cosine $1.000000$} \\
Gradient at the final position, thesis energy & exactly zero (the thesis's own theorem) & $\bm{0.000\mathrm{e}{+}00}$ \\
\bottomrule
\end{tabular}
\caption{Verification of the MuCoLa energy implementation.}
\end{table}

The third row reproduces the thesis's final-position theorem inside this code, and the
second row is precisely why that theorem does not apply to MuCoLa.

\subsection{The two energies head to head}

Identical code, optimizer, projection, initialization, step count and constraint weight.
The only difference is how the target token enters the language-model term.

\begin{table}[htbp]
\centering
\begin{tabular}{lrr}
\toprule
Variant & Steering-head acc. & Mean final LM log-prob \\
\midrule
\texttt{mucola\_faithful} & $40.0\%$ & $\bm{-38.5}$ \\
\texttt{cls\_thesis} & $53.3\%$ & $\bm{-353.9}$ \\
\bottomrule
\end{tabular}
\caption{Smoke scale, $n=15$, 60 steps. The steering-head column is deliberately
included; see the warning below.}
\end{table}

First generated continuation from each, same prompt and seed:

\begin{quote}
\texttt{mucola\_faithful}: ``Once upon a time, there was a king who ruled over a large
land. His subjects were very happy, and he''

\texttt{cls\_thesis}: ``Once upon a time environmentalCNN Unicornrosso PE \#\#YD vitality
ped Eh Karn bladesActionCodeStreamerBot Rothido Osborne''
\end{quote}

MuCoLa's energy produces fluent, coherent English. The energy this repository implements
produces word salad at a sequence log-probability roughly an order of magnitude worse.

\paragraph{The inverted steering number, and why a held-out judge is mandatory.}
\texttt{cls\_thesis} scores \emph{higher} on the steering classifier while producing
gibberish. A pooled sentiment head over hidden states is trivially satisfiable by
incoherent token soup, so an adherence number reported without a fluency number is
worthless. Every adherence figure in this programme is paired with an external fluency
figure for exactly this reason, and the classifier that steers is never the classifier
that scores.

\subsection{This confirms the submitted thesis rather than contradicting it}

\file{chapters/06\_discussion.tex} line 50 already states that the thesis ``refutes the
premise for a sampler that differentiates the input embedding of a likelihood whose
target token enters as a discrete index, and not for gradient-guided discrete sampling
as such'', and describes MuCoLa's self term correctly. The correction was applied before
submission. This experiment turns an argument made from the papers' text into a
measurement: implementing MuCoLa's energy in this repository's own code makes the
continuous sampler work, and implementing the repository's energy makes it fail, on the
same task with everything else held fixed.

\section{Study C: how should the constraint enter?}

\subsection{Shortlist coverage, reported before any adherence number}

The diagnostic: at each step, find the candidate the steering classifier most prefers
over the whole vocabulary, then ask where it falls in the likelihood ranking and whether
it lies inside the shortlist. Source
\file{results/ctg/phase2/Ccov.k\{16,64,256\}.json}, $n=30$ each.

\begin{table}[htbp]
\centering
\begin{tabular}{lrrrrr}
\toprule
$k$ & Coverage & Chance level & Ratio & Median likelihood rank & Linearization $\rho$ \\
\midrule
16  & $\bm{0.2\%}$ & $0.032\%$ & ${\sim}6\times$ & $29{,}928 / 50{,}257$ & $+0.019$ \\
64  & $\bm{0.8\%}$ & $0.127\%$ & ${\sim}6\times$ & $34{,}227 / 50{,}257$ & $+0.040$ \\
256 & $\bm{1.7\%}$ & $0.509\%$ & ${\sim}3\times$ & $29{,}866 / 50{,}257$ & $-0.008$ \\
\bottomrule
\end{tabular}
\caption{Shortlist coverage and the constraint-side linearization correlation.}
\end{table}

\paragraph{Finding 1: the shortlist almost never contains what the constraint wants.}
The median likelihood rank of the constraint's preferred token is around $30{,}000$ out
of $50{,}257$, the middle of the vocabulary: the two rankings are close to independent.
Coverage is three to six times chance, so the association is positive and real, but it
is still under $2$ percent even at $k=256$. Widening the shortlist sixteen-fold moves
coverage from $0.2$ to $1.7$ percent.

The consequence was pre-registered: any weakness of an exact-rescoring arm is a
\emph{proposal-support limitation}, not evidence about rescoring as a mechanism, because
such an arm can only rescore candidates the shortlist contains.

\paragraph{Finding 2: the constraint-side linearization carries no usable signal.}
The rank correlation between the classifier's first-order surrogate and the true change
in the classifier's score over the same candidates is $+0.019$, $+0.040$ and $-0.008$,
indistinguishable from zero.

\begin{table}[htbp]
\centering
\begin{tabular}{lr}
\toprule
Surrogate & Spearman against the true change \\
\midrule
Likelihood, token-indicator, admissible distances & $+0.60$ to $+0.73$ \\
Likelihood, input-embedding gradient & ${\sim}0.0$ \\
\textbf{Constraint, first-order classifier gradient} & $\bm{+0.02}$ \textbf{to} $\bm{+0.06}$ \\
\bottomrule
\end{tabular}
\caption{The classifier gradient is in the same condition as the input-embedding
gradient the thesis's null is about.}
\end{table}

\paragraph{A qualification, not a contradiction.}
The repository records the concern-11 reading as ``the constraint gradient's direction
carries signal, the LM gradient's does not'', resting on a paired contrast of roughly
$+27$ to $+37$ points on the continuation setup. That contrast is not re-run here and is
not disputed. What is added is that whatever signal the constraint direction carries is
not carried by the accuracy of its first-order expansion. A plausible reconciliation,
offered as a hypothesis rather than a result: a classifier gradient can point usefully in
the aggregate over many small continuous steps, which is what the continuous sampler
takes, while being a poor predictor of the effect of any single discrete token
substitution, which is what this measurement scores.

\paragraph{A tautological number, flagged so it is never misread.}
The mixture-shortlist arm reports $100.0\%$ coverage by construction, because it builds
half its shortlist from the constraint's own ranking. It is reported for completeness and
must never be presented as that arm achieving coverage.

\section{Study A on the constrained-generation task}

Partial at the time of writing; \texttt{exact\_k} and \texttt{uniform} were still
running. $n=60$ sequences, span $20$, $600$ steps, $k=64$, constraint in the accept step
only.

\begin{table}[htbp]
\centering
\begin{tabular}{lrrrrr}
\toprule
Arm & Steer-head acc. & Held-out adherence & Accept \% & fwd-eq/seq & GPU s \\
\midrule
\texttt{self}    & $36.7\%$ & $46.7\% \pm 12.6$ & 10.3 & $\bm{635}$ & $\bm{2{,}982}$ \\
\texttt{onehot}  & $41.7\%$ & $63.3\% \pm 12.2$ & 10.3 & $4{,}133$ & $8{,}047$ \\
\texttt{embgrad} & $38.3\%$ & $56.7\% \pm 12.5$ & 12.9 & $4{,}234$ & $10{,}712$ \\
\bottomrule
\end{tabular}
\caption{Held-out adherence is scored by an off-the-shelf sentiment classifier with no
connection to this project.}
\end{table}

The three adherence figures have $95\%$ intervals of roughly $\pm 12.5$ points on $n=60$
and overlap heavily, so the apparent ordering is \emph{not} a separation. The better
powered paired test, on the same prompt and seed across arms, puts both differences on
zero: \texttt{onehot} $-$ \texttt{self} gives $+0.138$ $[-0.186,+0.449]$ on the steering
log-probability, and \texttt{embgrad} $-$ \texttt{self} gives $+0.034$ $[-0.255,+0.320]$.

Consistent with the recovery result: the self term matches the full surrogate at $6.5$
times less compute. Q1's answer does not change when the task changes from repairing one
token to conjuring twenty.

\paragraph{One result that goes the other way.}
\texttt{embgrad} attains a higher final LM log-probability, $+12.52$ nats with a CI
excluding zero. This is not evidence of directional signal: every arm uses the identical
exact accept step and therefore targets the identical distribution, so a difference in
final energy is a difference in \emph{mixing}, and \texttt{embgrad}'s acceptance rate is
indeed higher. It is recorded rather than filtered out.

\section{Incidents, defects and harness lessons}

\subsection{Out-of-memory in the nearest-embedding projection}

Six jobs failed with CUDA out-of-memory at \file{core/prep.py} line 41. The shared
projection helper forms the difference tensor explicitly, of shape
$(M, V, D)$. At one masked position that is $0.26$ GB, which is why the entire
$145$-configuration recovery grid never touched it; at the twenty positions of the
continuation task it is $4.79$ GB in a single allocation, every step.

The fix is a local low-memory projection using $\lVert \vs - \ve \rVert^2 =
\lVert \vs\rVert^2 + \lVert \ve\rVert^2 - 2\vs^\top\ve$ and dropping the term constant in
the candidate. \file{core/prep.py} is deliberately not changed, because it is on the
archived path. Equivalence was verified on CPU over $120$ projections in three regimes:
\emph{zero mismatches}.

\subsection{A constraint-weight sweep that ran backwards}

The continuous-optimizer arms initially showed the steering score getting \emph{worse} as
the constraint weight increased, $-1.240 \to -1.268 \to -1.614$ at gains $0.5$, $2.0$,
$4.0$, with the generated text identical across gains. The cause is plain gradient ascent
without normalization: the gradient norm scales with the weights, so the step length
scales with the constraint weight and the iterate overshoots at the schedule's initial
step size, projecting to unrelated tokens.

The fix is to unit-norm the combined gradient per position before stepping. This leaves
the direction, and hence the weighting between the two terms, untouched, and makes the
step length independent of the weight. With it the sweep behaves monotonically
($73.3\% \to 80.0\% \to 80.0\%$ at gains $0.5$, $4.0$, $8.0$). This matters for fairness:
an unnormalized sweep would have understated MuCoLa's arm, which is the arm this
programme is at pains to represent fairly.

\subsection{Lessons that generalize}

\begin{itemize}
\item A single-mask diagnostic task hides a class of scaling bug. Two separate defects
  here were invisible at one masked position and appeared immediately at twenty.
\item Reversibility failures are silent in the worst way: acceptance collapses to exactly
  zero while the proposal itself is perfectly healthy. The diagnosis required measuring
  the proposal mass and the energy gain of the refused move.
\item The manifest is a priority list and can be reordered live; workers re-read it each
  pass. A gating experiment stuck behind a long family is a scheduling problem, not a
  compute problem.
\item Worker counts should be set by measurement. Two workers per card left the GPUs at
  $15$ to $43$ percent utilization on these chains; four per card reached $100$ percent.
\end{itemize}

\section{What was not done, stated rather than omitted}

\begin{itemize}
\item \textbf{Study B} (proposal sources including RoBERTa and SEDD on the constrained
  task) was gated on the Study C mechanism probe and was not launched. The harness
  supports both proposals and both were validated in the recovery setting by the earlier
  revision work.
\item \textbf{Study C's full mechanism comparison} at scale was in flight; only the
  coverage half is reported above as complete.
\item \textbf{Study D's} full adherence-fluency curve at matched compute was in flight.
\item \textbf{Study E}, diversity at matched adherence against a compute-matched
  best-of-$N$ optimizer, was not reached. The metrics are implemented and validated
  (distinct-1 $0.929$, distinct-3 $1.00$, self-BLEU $0.000$, semantic spread $0.778$ on
  one arm) but the controlled comparison is the part that matters, so no claim about
  posterior coverage is made on this evidence.
\end{itemize}

\section{Summary for the defense}

\begin{enumerate}
\item The constructive half of the thesis needs no gradient at all. The self term, which
  every autoregressive model already exposes, gets the whole effect at a quarter of the
  cost and zero backward passes.
\item The $40$ percent headline is an $n=50$ cell that reproduces exactly; at $n=200$ it
  is $26.5$ percent, so the honest claim is zero to roughly a quarter recovered, below
  SEDD and RoBERTa rather than at parity with them. The zero-versus-nonzero contrast,
  which is what the thesis rests on, is unaffected and better supported.
\item The scope restriction the thesis states about MuCoLa is correct, and can now be
  demonstrated rather than argued: the same code with MuCoLa's energy writes fluent
  English and with the thesis's energy writes word salad.
\item A shortlisted sampler needs a uniform-on-the-same-shortlist floor before any
  surrogate claim can be read, and that floor was never measured in the thesis.
\item Gradient-based constraint guidance is weak here for a measurable reason: the
  classifier's first-order expansion has a rank correlation near zero with the change it
  is expanding, and a likelihood-ranked shortlist almost never contains the token the
  constraint wants.
\end{enumerate}

\end{document}
